%============================================================
% Chapter 21: Offline and Parallel Binary Search
% STATUS: PLANNED  (generated by tools/outline.py; replace this file when the chapter is written)
%============================================================
\chapter{Offline and Parallel Binary Search}
\label{ch:offline}

\begin{tcolorbox}[enhanced, colback=white, colframe=secondaryblue, boxrule=1pt, arc=2mm]
\textbf{\large\textcolor{primaryblue}{Many queries, each needing its own binary search over the same process: run all of them together, round by round.}}
\end{tcolorbox}

\chapterfacts{2100}{2900}{3}{6}{30}

\begin{plannedbox}[Planned Contents --- this chapter is not written yet]
\textbf{Part A --- Lessons}
\begin{enumerate}[leftmargin=*, label=\textbf{A\arabic*.}, itemsep=1pt]
    \item The batched form: every query keeps $[lo, hi]$; each round replays the process once
    \item $\bigO{(n + q)\log V}$ rounds and the data structure per round
    \item Typical processes: adding updates in time order, adding edges by weight
    \item The recursive (divide and conquer on time) form is the recursion book (Ch 29)
    \item When a persistent structure or an offline sweep is simpler
\end{enumerate}
\textbf{Part B --- Patterns} (with candidate problems, checked when the chapter is written)
\begin{enumerate}[leftmargin=*, label=\textbf{P\arabic*.}, itemsep=1pt]
    \item \textbf{Batched binary search}
    \item \textbf{Updates in time order}
    \item \textbf{Thresholds on weights}
\end{enumerate}
\end{plannedbox}
